\documentclass[11pt,a4paper]{article}

\usepackage[a4paper,margin=2.2cm]{geometry}
\usepackage[indonesian]{babel}
\usepackage[T1]{fontenc}
\usepackage[utf8]{inputenc}
\usepackage{hyperref}
\usepackage{enumitem}
\usepackage{longtable}
\usepackage{array}
\usepackage{booktabs}
\usepackage{xcolor}

\hypersetup{
  colorlinks=true,
  linkcolor=black,
  urlcolor=blue,
  citecolor=black
}

\title{Daftar Tautan Unduhan Jurnal\\Proposal Skripsi FHE Backend}
\author{Muhammad Ulinnuha\\23/516588/PA/22101}
\date{\today}

\begin{document}
\maketitle

\noindent
Dokumen ini merangkum seluruh referensi yang dipakai pada proposal, lengkap dengan tautan unduhan (klik di PDF viewer yang mendukung hyperlink) dan lokasi file lokal di folder \texttt{refs/} bila sudah tersedia. Gunakan dokumen ini untuk mengisi arsip bacaan atau mengunduh ulang paper yang hilang.

\section*{Keterangan}
\begin{itemize}[leftmargin=1.2em]
  \item \textbf{OA} = open access / PDF gratis tanpa login (pada saat dicek).
  \item \textbf{DOI} = laman penerbit; kadang memerlukan akses kampus/IEEE/Springer.
  \item File lokal mengacu ke folder \texttt{refs/} di repositori proposal.
\end{itemize}

\section{Referensi utama proposal (12 entri)}

\begin{enumerate}[leftmargin=1.6em,itemsep=0.85em]
  \item \textbf{Devi \& Jayasri (2025)} --- Privacy Preserving Healthcare Based on Secure Mathematical Foundations of FHE.\\
  DOI: \url{https://doi.org/10.1109/ESCI63694.2025.10988128}\\
  Lokal: \texttt{refs/01\_Devi\_2025\_Privacy\_Healthcare\_FHE.pdf}

  \item \textbf{Peker \& Raj (2025)} --- Homomorphic Encryption for Confidential Statistical Computation.\\
  PDF MDPI (OA): \url{https://www.mdpi.com/2624-800X/6/1/4/pdf}\\
  HTML: \url{https://www.mdpi.com/2624-800X/6/1/4}\\
  Lokal: \texttt{refs/02\_Peker\_2025\_Confidential\_Statistical\_FHE.pdf}

  \item \textbf{Epishkina \& Ermakov (2025)} --- Homomorphic Encryption for Data Protection in Cloud Computing.\\
  DOI: \url{https://doi.org/10.1007/s11416-025-00552-6}\\
  Lokal: \texttt{refs/03\_Epishkina\_2025\_HE\_Cloud\_Computing.pdf}

  \item \textbf{Vinko et al.\ (2026)} --- Secure-by-Design FHE for Proficiency Testing / ILC.\\
  PDF MDPI (OA): \url{https://www.mdpi.com/2673-4001/7/1/14/pdf}\\
  HTML: \url{https://www.mdpi.com/2673-4001/7/1/14}\\
  Lokal: \texttt{refs/04\_Vinko\_2026\_Secure\_by\_Design\_FHE.pdf}

  \item \textbf{Wu et al.\ (2025)} --- Homomorphic Encryption for ML with CKKS: A Survey.\\
  Penerbit: \url{https://www.techscience.com/cmc/v85n1/62800}\\
  DOI: \url{https://doi.org/10.32604/cmc.2025.064346}\\
  Lokal: \texttt{refs/05\_Wu\_2025\_CKKS\_ML\_Survey.pdf}

  \item \textbf{Li et al.\ (2026)} --- Survey of Optimization Techniques for Bootstrapping Algorithms in FHE.\\
  DOI / Springer (OA): \url{https://doi.org/10.1186/s42400-026-00571-w}\\
  Lokal: \texttt{refs/06\_Li\_2026\_Bootstrapping\_Survey.pdf}

  \item \textbf{Fan et al.\ (2026)} --- Understanding and Boosting FHE Applications on GPU.\\
  DOI / Springer (OA): \url{https://doi.org/10.1186/s42400-025-00482-2}\\
  Lokal: \texttt{refs/07\_Fan\_2026\_FHE\_on\_GPU.pdf}

  \item \textbf{Xu et al.\ (2026)} --- HERA: Bandwidth-efficient FHE Accelerator on HBM FPGA.\\
  DOI ACM: \url{https://doi.org/10.1145/3748173.3779201}\\
  Lokal: \texttt{refs/08\_Xu\_2026\_HERA\_FPGA.pdf}

  \item \textbf{Lee et al.\ (2025)} --- Homomorphic Evaluation Cluster Architecture for FHE.\\
  DOI IEEE Open: \url{https://doi.org/10.1109/OJCAS.2025.3568058}\\
  Lokal: \texttt{refs/09\_Lee\_2025\_HE\_Cluster\_Architecture.pdf}

  \item \textbf{Al Badawi et al.\ (2022)} --- OpenFHE: Open-Source Fully Homomorphic Encryption Library.\\
  PDF ePrint (OA, wajib untuk stack prototipe): \url{https://eprint.iacr.org/2022/915.pdf}\\
  Halaman: \url{https://eprint.iacr.org/2022/915}\\
  Lokal (jika sudah diunduh): \texttt{refs/10\_Badawi\_2022\_OpenFHE.pdf}

  \item \textbf{Brynds et al.\ (2025)} --- CryptOracle: Modular Framework to Characterize FHE (OpenFHE/CKKS).\\
  PDF arXiv (OA): \url{https://arxiv.org/pdf/2510.03565.pdf}\\
  Abs: \url{https://arxiv.org/abs/2510.03565}\\
  Lokal: \texttt{refs/11\_Brynds\_2025\_CryptOracle.pdf}

  \item \textbf{Benaissa et al.\ (2021)} --- TenSEAL: Encrypted Tensor Operations Using HE (jalur cadangan SEAL).\\
  PDF arXiv (OA): \url{https://arxiv.org/pdf/2104.03152.pdf}\\
  Abs: \url{https://arxiv.org/abs/2104.03152}\\
  Lokal: \texttt{refs/12\_Benaissa\_2021\_TenSEAL.pdf}
\end{enumerate}

\section{Ringkasan cepat (salin ke browser)}

\begin{verbatim}
https://www.mdpi.com/2624-800X/6/1/4/pdf
https://www.mdpi.com/2673-4001/7/1/14/pdf
https://doi.org/10.1186/s42400-026-00571-w
https://doi.org/10.1186/s42400-025-00482-2
https://eprint.iacr.org/2022/915.pdf
https://arxiv.org/pdf/2510.03565.pdf
https://arxiv.org/pdf/2104.03152.pdf
https://doi.org/10.1109/ESCI63694.2025.10988128
https://doi.org/10.1007/s11416-025-00552-6
https://doi.org/10.32604/cmc.2025.064346
https://doi.org/10.1145/3748173.3779201
https://doi.org/10.1109/OJCAS.2025.3568058
\end{verbatim}

\section{Catatan unduhan}
\begin{itemize}[leftmargin=1.2em]
  \item Paper arXiv dan IACR ePrint biasanya paling mudah diunduh tanpa login.
  \item MDPI dan Springer Open Access juga menyediakan tombol PDF langsung.
  \item IEEE/ACM kadang memerlukan jaringan kampus atau akun institusi; salinan lokal di \texttt{refs/} sudah disiapkan bila file sumber ada di folder proposal.
  \item Setelah mengunduh OpenFHE (no.\ 10), simpan sebagai \texttt{refs/10\_Badawi\_2022\_OpenFHE.pdf} agar penamaan konsisten.
\end{itemize}

\section{Apakah perlu scaffold backend sekarang?}
Pada tahap \textbf{proposal}, scaffold OpenFHE/FastAPI \textbf{belum wajib}. Proposal sudah memuat stack, endpoint, beban kerja, dan kriteria uji. Scaffold menjadi relevan setelah proposal disetujui, yaitu saat masuk implementasi skripsi (tonggak 1--4 di BAB IV). Memulai scaffold terlalu dini berisiko mengubah desain sebelum umpan balik pembimbing masuk.

\end{document}
